\chapter{Efficiency Tips: Work Like a Pro}

This chapter collects operating habits that make a big difference once you know them. Each one is short,
but together they can double your modelling speed.

\section{Using gestures efficiently}

\begin{center}
\small
\begin{tabularx}{\linewidth}{@{}lX@{}}
\toprule
\textbf{Situation} & \textbf{Tip} \\
\midrule
Finger covers the target & Long-press to enter the \textbf{Precision pointer}: a crosshair appears 60 dp above your finger, and lifting the finger completes the pick \\
Cannot hit a small element & Zoom in with two fingers first, then tap; picking has a priority order (vertex 12 dp → edge 10 dp → face) \\
Need to rotate the sketch view & In a sketch one finger draws, and it takes \textbf{three fingers} to rotate; alternatively tap the navigation cube \\
Clear the selection by accident & A double tap means ``clear selection'', so avoid double taps while making precise selections \\
Want a right-click & \textbf{Two-finger tap} on the view opens a menu equivalent to the desktop right-click menu \\
Select many objects at once & \textbf{Two-finger long-press} starts box selection; drag a rectangle, and use the selection filter to box only faces or only objects \\
Leave the current tool & The toolstrip at the bottom-right has ✕ (the same as Esc); or press the back button \\
\bottomrule
\end{tabularx}
\end{center}

\section{The back button is a universal Esc}

Remember this one rule: \textbf{the back button closes one layer at a time}, starting from the innermost---numpad,
text input, picker, menu, tool panel, Draft, assembly joint, task panel, feature dialog, sketch. Once everything
is closed, the app goes to the background instead of quitting. So whenever you are stuck in a mode, pressing back
repeatedly always gets you back to normal.

\section{Selection and selection filters}

\begin{itemize}
  \item A tap \textbf{adds to the selection}, and tapping the same element again removes it; the ✕ on the selection bar removes one item; a double tap clears everything.
  \item To select only faces or only edges, use the selection filters in \menupath{Tools} (or commands such as \cmd{Part\_FaceSelection}):
        a ``faces only'' strip appears in the bottom-right corner of the view, and you tap ✗ to clear the filter.
  \item Selection history: \menupath{Edit → Selection back/forward} (\cmd{Std\_SelBack}/\cmd{Std\_SelForward}) moves back and forth between selection sets.
\end{itemize}

\section{Use expressions instead of numbers}

Every numeric input accepts an expression: type it straight into the top field of the numpad, or start it with \file{=}. For example:

\begin{lstlisting}
=60 mm / 2          // 30 mm
=Sheet.Thickness*2  // refers to a spreadsheet alias (Chapter 11)
=10 mm + 5 mm       // arithmetic with units
\end{lstlisting}

The benefit: change one variable and every dimension that references it changes as well. In complex models,
put the key dimensions in a spreadsheet (the Variables tab) and reference them with expressions---that is the
most trouble-free approach.

\section{Replay the modelling history with Tip}

A body's \ui{Tip} marker decides how far the model is computed. Long-press a feature in the tree and choose
\ui{Roll tip to here}; the model rolls back to the state at that step, so you can inspect intermediate results
and make changes to early features. Roll forward to the last feature again to restore the full model.

Together with \ui{Suppress} you can leave a feature out temporarily and compare the result with and without it;
errors caused by that feature disappear along with it.

\section{Tool search and recent tools}

\begin{itemize}
  \item \ui{Search all tools} at the top of the tool panel is a global search: type the command's Chinese name,
        its English name (pad), or its command ID (PartDesign\_Pad) and it matches;
  \item ``Recent tools'' on the right of the toolstrip keeps 6 commands, so the ones you use repeatedly stay at hand;
  \item Long-pressing a button in the tool panel does not pop up a description; the name shown is the FreeCAD command name,
        which you can look up in Appendix A and in the online documentation.
\end{itemize}

\section{Keep your most-used panels in the Variables tab}

A spreadsheet is more than arithmetic: write your design parameters (length, width, height, hole diameter,
pattern count) into one table and reference it everywhere while modelling. Change one cell and the whole
model updates. This is the cockpit of parametric design.

\section{Performance and smoothness}

\begin{itemize}
  \item \textbf{Hide what you are not looking at}: in a large assembly, hide unrelated components (long-press in the tree → Hide) and the view redraws noticeably faster;
  \item \textbf{Keep fewer dialogs open at once}: feature dialogs preview in real time and keep recomputing on complex models;
  \item \textbf{Model in parts}: one body with hundreds of features is worse than several bodies brought together in an assembly;
  \item \textbf{FEM meshes}: on mobile, run a coarse mesh first and refine once the boundary conditions are confirmed; plug in the charger and keep the screen awake while solving;
  \item \textbf{Storage}: the runtime keeps about 1 GB resident, and when free space drops below 2 GB the system may kill background apps.
\end{itemize}

\section{Hardware keyboards and mice}

With a keyboard or mouse connected to the tablet, you can type straight into input fields and the layout is not
rebuilt (plugging in a keyboard does not interrupt what you are doing). The left mouse button acts as a tap, the right
button as a context menu, and the wheel as zoom (the same as a two-finger pinch). Multi-finger trackpad gestures are
handled by the system.

\section{Autosave and recovery}

MobileCAD writes the current document to a recovery copy every 60 seconds and every time it goes to the background.
If the app crashes, the next start opens the recovery copy automatically and shows ``Recovered unsaved work from HH:mm'';
save it to a real file at once. The recovery copy is deleted after a normal save.
